\documentclass[conference]{IEEEtran}
\IEEEoverridecommandlockouts

% margins: 2.5 cm top and bottom; 3 cm left and right
\usepackage[
a4paper,
vmargin=2cm,
hmargin=1.5cm
]{geometry}
%\usepackage{cite}
% Paragraph Spacing and Line Spacing: Narrow (6 pt / 1.15 spacing) or Moderate (6 pt / 1.5 spacing)
%\renewcommand{\baselinestretch}{1.15}
%\parskip=6pt
\usepackage{enumitem}

% Color settings for cover and code listings 
\usepackage[table]{xcolor}
\definecolor{azulUC3M}{RGB}{0,0,102}
\definecolor{gray97}{gray}{.97}
\definecolor{gray75}{gray}{.75}
\definecolor{gray45}{gray}{.45}

% PDF/A -- Important for its inclusion in e-Archive. PDF/A is the optimal format for preservation and for the generation of metadata: http://uc3m.libguides.com/ld.php?content_id=31389625. 

% In the template we include the file OUTPUT.XMPDATA. You can download that file and include the metadata that will be incorporated into the PDF file when you compile the memoria.tex file. Then upload it back to your project.  
\usepackage[a-1b]{pdfx}

% LINKS
\usepackage{hyperref}
\hypersetup{colorlinks=true,
	linkcolor=black, % links to parts of the document (e.g. index) in black
	urlcolor=blue} % links to resources outside the document in blue

% MATH EXPRESSIONS
\usepackage{amsmath,amssymb,amsfonts,amsthm}

% Character encoding
\usepackage{txfonts} 
\usepackage[T1]{fontenc}
\usepackage[utf8]{inputenc}

% English settings
\usepackage[english]{babel} 
\usepackage[babel, english=american]{csquotes}
\AtBeginEnvironment{quote}{\small}

% Footer settings
\usepackage{fancyhdr}
\pagestyle{fancy}
\fancyhf{}
\renewcommand{\headrulewidth}{0pt}
\rfoot{\thepage}
\fancypagestyle{plain}{\pagestyle{fancy}}


% Tables and figures settings
\usepackage{multirow} % combine cells 
\usepackage{caption} % customize the title of tables and figures
\DeclareCaptionOption{compatibility}{false}
\usepackage{floatrow} % we use this package and its \ ttabbox and \ ffigbox macros to align the table and figure names according to the defined style.
\usepackage{array} % with this package we can define in the following line a new type of column for tables: custom width and centered content
\newcolumntype{P}[1]{>{\centering\arraybackslash}p{#1}}
\DeclareCaptionFormat{upper}{#1#2\uppercase{#3}\par}
\usepackage{graphicx}
\graphicspath{{imagenes/}} % Images folder

% Table layout for engineering
\captionsetup*[table]{
	format=hang,
	name=Table,
	justification=centering,
	labelsep=period,
	width=.55\linewidth,
	labelfont=small,
	font=small
}

% Figures layout for engineering
\captionsetup[figure]{
	format=hang,
	name=Fig.,
        justification=centering,
	singlelinecheck=off,
	labelsep=period,
	labelfont=small,
	font=small		
}

% FOOTNOTES
\usepackage{chngcntr} % continuous numbering of footnotes


% CODE LISTINGS 
% support and styling for listings. More information in  https://es.wikibooks.org/wiki/Manual_de_LaTeX/Listados_de_código/Listados_con_listings
\usepackage{listings}

% Custom listing
\lstdefinestyle{estilo}{ frame=Ltb,
	framerule=0pt,
	aboveskip=0.5cm,
	framextopmargin=3pt,
	framexbottommargin=3pt,
	framexleftmargin=0.4cm,
	framesep=0pt,
	rulesep=.4pt,
	backgroundcolor=\color{gray97},
	rulesepcolor=\color{black},
	%
	basicstyle=\ttfamily\footnotesize,
	keywordstyle=\bfseries,
	stringstyle=\ttfamily,
	showstringspaces = false,
	commentstyle=\color{gray45},     
	%
	numbers=left,
	numbersep=15pt,
	numberstyle=\tiny,
	numberfirstline = false,
	breaklines=true,
	xleftmargin=\parindent
}

\captionsetup*[lstlisting]{font=small, labelsep=period}
 
\lstset{style=estilo}
\renewcommand{\lstlistingname}{\uppercase{Código}}


% REFERENCES 
% IEEE bibliography setup
\usepackage[backend=biber, style=ieee, isbn=false,sortcites, maxbibnames=6, minbibnames=1]{biblatex} % Setting for IEEE citation style, recommended for engineering. "maxbibnames" indicates that from 6 authors truncate the list in the first one (minbibnames) and add "et al." as used in the IEEE style.

% Añadimos las siguientes indicaciones para mejorar la adaptación del estilos en español
\DefineBibliographyStrings{spanish}{%
	andothers = {et\addabbrvspace al\adddot}
}
\DefineBibliographyStrings{spanish}{
	url = {\adddot\space[En línea]\adddot\space Disponible en:}
}
\DefineBibliographyStrings{spanish}{
	urlseen = {Acceso:}
}
\DefineBibliographyStrings{spanish}{
	pages = {pp\adddot},
	page = {p.\adddot}
}

\addbibresource{references.bib} % llama al archivo referencias.bib en el que deberá estar la bibliografía utilizada



%-------------
%	DOCUMENT
%-------------

\begin{document}
	
%----------
%	COVER
%----------	
\begin{titlepage}
	\begin{sffamily}
	\color{azulUC3M}
	\begin{center}
		\begin{figure}[H] % UC3M Logo
			\makebox[\textwidth][c]{\includegraphics[width=16cm]{images/logo_UC3M.png}}
		\end{figure}
		\vspace{2.5cm}
		\begin{Large}
			Master Degree in Cybersecurity\\			
			 2025-2026\\ % Academic year
			\vspace{2cm}		
			\textsl{Master Thesis}
			\bigskip
			
		\end{Large}
		 	{\Huge ``Security Analysis of Modern Tor V3 Hidden Service Directories"}\\
		 	\vspace*{0.5cm}
	 		\rule{10.5cm}{0.1mm}\\
			\vspace*{2cm}
			{\LARGE Alejandro García Berrocal}\\ 
			\vspace*{3cm}
		\begin{Large}
			Marcelo Gabriel Bagnulo Braun  \\
			
			Madrid, 14 Sep 2026\\
		\end{Large}
	\end{center}
	\vfill
	\color{black}
	\fbox{
	\begin{minipage}{\linewidth}
    	\textbf{AVOID PLAGIARISM}\\
    	\footnotesize{The University uses the \textbf{Turnitin Feedback Studio} for the delivery of student work. This program compares the originality of the work delivered by each student with millions of electronic resources and detects those parts of the text that are copied and pasted. Plagiarizing in a TFM is considered a  \textbf{\underline{Serious Misconduct}}, and may result in permanent expulsion from the University.}\end{minipage}}

	% IF OUR WORK IS TO BE PUBLISHED UNDER A CREATIVE COMMONS LICENSE, INCLUDE THESE LINES. IS THE RECOMMENDED OPTION.
	\noindent\includegraphics[width=4.2cm]{images/creativecommons.png}\\ % Creative Commons Logo
    \footnotesize{This work is licensed under Creative Commons \textbf{Attribution – Non Commercial – Non Derivatives}}
	
	\end{sffamily}
\end{titlepage}



\title{Security Analysis of Modern Tor V3 Hidden Service Directories}

\author{\IEEEauthorblockN{Alejandro García Berrocal}
\IEEEauthorblockA{\textit{Student}, \textit{Carlos III University}, Madrid, Spain,\\
100451059@alumnos.uc3m.es}
}

\maketitle

\begin{abstract}
TODO
\end{abstract}

\begin{IEEEkeywords} %TODO: Modify these keywords
Tor, Hidden Services, HSDirs, Cybersecurity, Privacy, Analysis, Churn, Sybil patterns
\end{IEEEkeywords}

\section{Introduction} %Refactor the initial writing style introduced in the introduction
The Tor Project, commonly known as Tor, it is, if not the most, preferred network overlay to provide 
an anonymous communication network. In 2003 the developers introduced an anonymity mechanism to protect
the service providers, onion services, on the Tor network, that was later deployed in 2004 as the version 2 (V2)
protocol of the onion services. %(TODO: USE REF 1 OF DOC)

Each onion service generates an address when started, a domain name, used by Tor clients to access the service.
To achieve a communication between a client and an onion service a two-way anonymous communication is established
over a 6-hop circuit, beaing each hop a Tor node. Nonetheless, Tor deprecates the version 2 protocol
in 2021 (TODO: USE REF 2 OF DOC), and transitions to a third version of onion services (V3) due to the 
different issues identified on the V2 protocol. Among these issues are included, the V2 protocol had unsecure cryptography
under today's technology capabilities (SHA1, DH \& RSA1024), it suffered from information leakeage to directory servers allowing attackers 
to harvest onion addresses by deploying their owned Hidden Service Directories, now onward referred as HSDirs, or performing enumeration attacks
towards other HSDirs on the Tor network. The newly deployed V3 protocol assesses these issues by including measurements to avoid a
direct exposure of onion addresses to HSDirs, introduced better cryptographic functions (SHA3, Ed25519 \& Curve25519), reduced the amount 
of information exposed to directory servers (HSDirs) and attack surface, improved security against impersonation, included an extensile of introduction 
and rendezvous points on the protocol, set offline keys (local) for onion services and other protection mechanisms. %(TODO: TOR DOC PAG 339)

To comprehend the current situation of HSDirs security, a set of onion services, consisting of personal onion services static
and ephemeral, and third party onion services, are put under a 30-days public data gathering of HSDirs hosting their 
onion services descriptors and network-wide HSDirs available on the network, to perform a comprehensive analysis to shed light on the current stability and security regarding
hosting HSDirs (Tor relays) on the Tor network. 

%TODO: Explanation of what it's been done and part of the results
\section{Tor's Onion Service V3 Protocol and Mechanism}

This section presents the Tor's onion service V3 protocol with an overview of the essential parties of the
mechanism, including the onion services address generation, publication and indexing, and connection establishment. %(TODO: USE TOR OFFICIAL DOCS AS REF)

\subsection{Mechanism and Protocol Overview}%TODO: Add references from Tor's official sources and infographics

The architectural ecosystem of the Tor version 3 (V3) onion service protocol implements a decentralized framework designed to facilitate mutual anonymized application-layer
communication for both the client and the service provider while masking the identities and locations of both endpoints. Different from centralized domain registries, the V3
architecture operates entirely within the Tor relay network, managing connections through a specialized set of relays avoiding direct network visilibility of both endpoints.
The public service address is splitted from their current IP routing path with a transient overlay and a multi-hop circuit. The framework relies essential in the interaction
of six operational entities distributed across the network:

\begin{enumerate}
    \item \textbf{Tor Client:} The end-user application seeking to establish an authenticated, secure connection to the onion service without disclosing its own identity. Commonly a local proxy dawmon that encapsulates the traffic, handles keys exchanges and establishes multi-hop virtual paths to mask the origins requests.
    \item \textbf{Onion Service:} The destination instance and isolated server hosting a service or application endpoints that listen uniquely to incoming traffic through a pre-arranged anonymous path while preserving the confidentility of its public IP address and network topology.
    \item \textbf{Directory Authorities:} A cluster of core trusted infrastructure server nodes that collaboratively manage and monitor the network consensus status, calculate the daily Shared Random Values (SRVs), and issue the the network consensus directory each hour required for layout lookups. 
    \item \textbf{Introduction Point - IPO:} A set of publicly accessible Tor relays arbitrarily selected by the onion service to maintain long-lived and persistent circuits. The operational duty of these nodes is to accept received incoming initialization requests from clients and proxying them to the service via pre-established circuits without knowing its actual location.
    \item \textbf{Hidden Service Directory - HSDir:} A dynamic subset of Tor relays that meet strict stability and uptime criteria. These nodes are authorized to store and serve signed cryptographic service descriptors within a decentralized and time dependent Distributed Hash Table (DHT) structure known as Hash Ring.
    \item \textbf{Rendezvous Point - RPO:} Ephemeral relays dynamically chosen by the client on the spot. These nodes act as the neutral nexus where the circuit branches from the client and service safely converge.
\end{enumerate}

%(INCLUDE IMAGE THAT EXPLAINS THE NETWORK)

The chronological lifecycle required to establish an end-to-end connection between the client and a V3 onion service transitions through a six cryptographic and operational stages:

\begin{enumerate}
    \item \textbf{IPO Setup:} Upon initialization, the V3 onion service generates its root identity key pairs, derives its cryptographic \texttt{.onion} address and establishes individual three-hop circuits to multiple chosen public relays to act as active IPOs.
    \item \textbf{Descriptor Publication:} The service generates an onion service descriptor containing its public verification keys and the network routing parameters of its IPOs. This document is cryptographically signed using an Ed25519 key-blinding mechanism that daily rotates, and uploaded to a default of eight designated HSDir positions derived methematically on the Hash Ring.
    \item \textbf{Discovery and Descriptor Retrieval:} The client to reach the target onion address queries the active network consensus along the global SRV to pinpoint and compute the exact HSDir relays responsible of hosting the service's current descriptor and downloading the authenticated record over an anonymous circuit.
    \item \textbf{Rendezvous Establishment:} The client designates an independent Tor relay from the active consensus pool to act as RPO, builds a separate three-hoop circuit to it and establish a single-use cryptographic secret denominated as rendezvous cookie.
    \item \textbf{Introduction Delivery:} The client dispatches an encrypted cell (\texttt{INTRODUCE1}), containing the RPO routing descriptors and the rendezvous cookie and encrypted using the service's public key, to one of the service's published IPOs, which blindly forwards it through its pre-established circuit to the hidden service.
    \item \textbf{Circuit Splicing:} Upon decrypting the cell, the onion service parses the IPO cell, establishes a newly three-hop circuit to the specified RPO and handovers a validation via a (\texttt{RENDEZVOUS1}) cell. The RPO then confirms the cookie matches and joints the two independent streams into a unified six-hop end-to-end circuit.
\end{enumerate}

\subsection{Address Generation}%TODO: Add references from Tor's official sources and infographics

The Tor version 3 (V3) onion routing introduces a more cryptographically and decentralized naming convention replacing the legacy 16-character version 2 (V2) address space which derivation introduced significant cryptographic
vulnerabilities including partial-match dictionary collisions, namespace harvesting and lack of support for forward-secure operations. The V3 protocol remediates these vulnerabilities by transitioning to Ed25519 elliptic curve
cryptography mapping the public key directly into the string encoding of the public identifier serving as self-authentication. This shift increases the \texttt{.onion} address length into a 56-character base32-encoded string,
eliminating central name authorities dependencies and preventing unauthorized directory harversting from mapping or enumerating hidden application endpoints while proving ownership.

The generation of a valid V3 onion address consists in the congregation of a 35-byte binary payload composed of three crucial fields, then serialized using Base32 encoding:

\begin{enumerate}
    \item \textbf{Master Public Key (32 bytes):} The Ed25519 master public key ($\mathcal{P}_{k}$) generated by the onion service during its initialization phase and root identity of the hidden service.
    \item \textbf{Verification Checksum (2 bytes):} A truncated security chunk ($\mathcal{C}_{s}$) calculated derived from cryptographic hash function to ensure string integrity and protect against typos or transmission distortions.
    \item \textbf{Protocol Version (1 byte):} A fixed administrative byte ($\mathcal{V}$) indicating the technical generation of the protocol explicitly configured to \texttt{0x03} for Tor version 3 (V3).
\end{enumerate}

To calculate the 2-byte verification checksum, the protocol combines a predefined static string \texttt{".onion checksum"} as a fixed salt, the 32-byte master public key, and the the protocol version byte into a unified byte stream. 
The stream is processed using the SHA3-256 hashing algorithm, which resulting digest is truncated to its first two bytes:
\begin{equation}
\mathcal{C}_{\text{s}} = \text{SHA3-256}\big(\textit{".onion checksum"} \parallel \mathcal{P}_{\text{k}} \parallel \mathcal{V}\big)\big[:2\big]
\end{equation}
Once the checksum ($\mathcal{C}_{\text{s}}$) is derived, it is appended to the master public key and the protocol version byte to construct the 35-byte vector. The resulting domain name is produced by converting the 
byte vector into an Base32 sequence with the domain suffix appended:
\begin{equation}
\mathcal{A}_{\text{onion}} = \text{Base32}\big(\mathcal{P}_{\text{k}} \parallel \mathcal{C}_{\text{s}} \parallel \mathcal{V}\big) \parallel \textit{".onion"}
\end{equation}
%TODO: Illustration

This infrastructure produces direct implications for network telemetry and monitoring. Because the resulting onion address wraps the raw master key ($\textit{P}_{k}$), any external system that captures the service's address
can extract the inherent cryptographic key. The master identity key acts as the fundamental parameter to calculate daily blinded keys, which govern the exact time dependent coordinates the service will occupy on the network's HRT
over the given lifecycle.

\subsection{Publication and Indexing}%TODO: Add references from Tor's official sources and infographics

To shield hidden services from targeted denial-of-service (DoS) vectors and metadata harversting by adversary directory servers, the Tor version 3 (V3) introduces a key-blinding subsystem along a consistent distributed hashing topology. 
Rather than uploading static routing records to fixed servers, V3 onion services routinely rotates ephemeral blinded public keys ($\textit{K}_{b}$) to dynamically compute a set of time dependent storage nodes within the HRT. Services
ensure confidentiality by serializing active routing vectors and cryptographic information in a structural document known as onion descriptor, protected with an architectural encryption to prevent intermediate nodes from harvesting onion addresses
or identify the services being hosted by it unless they already posses the corresponding credentials.

The derivation of blinded credentials, spatial positioning and temporal rotation descriptors accross the global HRT shifts are governed by two synchronized parameters daily updated within the network consensus, the global SRV and the time period, or temporal epoch. 
The consensus cluster of directory authorities produce a newly generated SRV every consecutive interval (epoch) defined by the consensus parameter \texttt{hsdir-interval}, which defaults to an operational duration of 1440 minutes (24 hours), at 00:00 UTC over hourly voting intervals:

\begin{enumerate}
    \item \textbf{Commit Phase (00:00 - 12:00 UTC):} Authorities include cryptographic commitments of privetly generated random numbers within their hourly consensus votes.
    \item \textbf{Reveal Phase (12:00 - 00:00 UTC):} Authorities expose their random values, validating them against the previous phase's commitments. At 00:00 UTC, the final unified SRV is compiled and appended to the live consensus.
\end{enumerate}

This shared secret acts a global seed that inforces pseudo-random reshuffling of the entire hash ring, preventing pre-calculation and maintance of malicious nodes adjacent to a target's publication space.

To publish its onion descriptor, the hidden service applies a key-blinding function mixed with the two current synchronized temporal variables to produce the temporal blinded public key $\textit{K}_{\text{b}}$. For each independent publication replica the service derives a distinct 256-bit
coordinate refer as the service index, where the number of replicas ($\textit{N}_{\text{r}}$) defaults to 2 replicas. The derivation of the service index for a given replica identifier ($\textit{r}$) is as follows:
\begin{equation}
\mathcal{HS}_{\text{index}} = \text{SHA3-256}\big(\textit{"store-at-idx"} \parallel \textit{K}_{\text{b}} \parallel \textit{r} \parallel \textit{L}_{\text{p}} \parallel \textit{N}_{\text{p}}\big)
\end{equation}
where $\textit{K}_{\text{blinded}}$ represents the blinded public key, $\textit{L}_{\text{p}}$ is the standardized length of the temporal epoch, and $\textit{N}_{\text{p}}$ is the sequential index of the current time period. 

Concurrently, every active relay in the consensus that satisfies specific uptime, stability, and flag requirements (\texttt{HSDir, Stable, Fast}) to operate as HSDir computes its own relative position on the ring, recalculated for each temporal epoch using the global SRV:
\begin{equation}
\mathcal{R}_{\text{index}} = \text{SHA3-256}\big(\textit{"node-idx"} \parallel \textit{ID} \parallel \text{SRV} \parallel \textit{N}_{\text{p}} \parallel \textit{L}_{\text{p}}\big)
\end{equation}
where $\textit{ID}$ denotes the cryptographic fingerprint of the relay.
%TODO: Illustration
Once both indeces are mapped onto the 256-bit circular index table of the HRT, the onion identifies the specific relays in charge of reduntdantly storing the current onion descriptor. The service selects the first distinct HSDir nodes ($\textit{k}$) whose positions on the HRT 
($\mathcal{R}_{\text{index}}$) immediately succeed the hidden service positions ($\mathcal{HS}_{\text{index}}$), where the pool number of nodes ($\textit{k}$) is given by the consensus parameter \texttt{hsdir-spread-store} explicitly set to 4. Thus, the service initiates a sequence
of encrypted descriptor uploads to these first 4 unique HSDir nodes indentified for each service index computed, resulting in the descriptor upload to total of 8 distinct HSDir nodes to host and serve the record during the designated time period.

Because these positions shift daily along the global SRV, this parameter along $\textit{N}_{\text{p}}$, $\textit{L}_{\text{p}}$ are publicly observable through the public network consensus files which serve to track a targeted \texttt{.onion} address with an independent system
to extract the onion's master key, compute $\textit{K}_{\text{b}}$, and determine the exact 8 HSDir relays responsible of hosting the onion's descriptor any given day. This mapping provides the basis to track HSDir node uptime distributions, assess runtime stability or operator entity 
cluster profiling via IP range correlation.

\subsection{Service Access} %TODO: Add references from Tor's official sources and infographics

Unlike the classical Internet, where a client establishes a direct transport layer socket to an IP address, the Tor version 3 (V3) requires both the client and server to construct separate three-hop virtual circuits which congregation to a mutually agreed intermediate node is referred as
rendezvous mechanism. This architecture ensures neither end can deduce the routing trajectory, geo-origin and underlying network topology of the opposite party. 

The service access process begins with the client parsing the 56-character \texttt{.onion} address to locate the master public identity key and resolves the inherent Ed25519 public key, which combined with the temporal epoch and SRV retrieved from the network status consensus, the client
cryptographically computes the target service $\textit{K}_{\text{b}}$ and derive the settled service indices mapped onto the HRT. The client then selects one of the four active HSDir relays appointed to the replica, constructs an anonymous three-hop circuit toward it, and transmits an authenticated
HTTP GET request to download the signed onion descriptor records.

Upon receiving the payload, the client parses the descriptor utilizing the blinded key structures to verify the descriptor's cryptographic signature and decrypt its inner layer. This extracts the array of active IPOs routing vectors along with their respective public keys and link specifiers, and the 
service's temporal encryption keys. The connection sequence then diverges into two parallel paths:

\begin{enumerate}
    \item \textbf{Rendezvous Establishment:} The client selects an external relay from the consensus to act as an isolated cryptographic bridge referred as RPO. A three-hop circuit is constructed toward this node and delivered with an \texttt{ESTABLISH\_RENDEZVOUS} cell, which provisions the relay with a secret and an unique single-use 20-byte rendezvous cookie. Upon acknowledging the setup, the RPO registers the cookie in its active volatile session cache and awaits an inbound circuit connection that presents the corresponding token.
    \item \textbf{Introduction Delivery:} Simultaneously, the client constructs a three-hop circuit to one of the public IPOs specified in the decrypted descriptor records. An \texttt{INTRODUCE1} authenticated cell is then transmitted to it, which envelops an encrypted message containing the routing specifiers and the matching 20-byte rendezvous cookie encrypted and safeguarded under the service's onion key against intermediate inspection.
\end{enumerate}

Upon receiving this cell, the IPO forwards the encrypted inner payload to the onion service via an \texttt{INTRODUCE2} cell. The hidden service then decrypts this instruction block using its blinded private key, validates its integrity, draws out the RPO's routoing parameters alongside the 20-byte rendezvous cookie
and establishes a separate three-hop circuit to the specified RPO relay by submitting a \texttt{RENDEZVOUS1} cell embedded with the matching rendezvous cookie. The RPO finally splices the incoming circuit into a unified six-hoop transparent pipeline facilitating end-to-end forward-secure data streams protected with 
symmetric encryption keys safeguarding against node sniffing or network-wide traffic correlation attack.

\section{Methodology}

To investigate and analyze the architectural security, runtime, stability, behaviour and Sybil deployments within Tor version 3 (V3) HSDirs, a non-intrusive and reliable tracking framework is designed. Rather than relying in on historical third-party directory archives which lack granularity, an automated and concurrent 
program was developed and executed during one-month time-period. The measurement infrastructure interacts directly with a local Tor instance via its asynchronous control port pipeline. The core approach of this methodology is the real-time interception of routing descriptors and propagation, track individual relay flag 
churn, state fluctuations, availability on the HRT and isolate structural network identifiers from the live network.

\subsection{Concurrent System Architecture and State Synchronization}

The measurement framework operates continuously over a set of decoupled multi-threaded structure to isolate telemetry vectors and maintain a minimal overload and performance of Tor's control port socket. The high throughput and directory update rates across the network would bring forth thread blockages and risk frame 
drops on the asynchronous controller stream in a single-threaded architecture. Consequently, the architecture separates its duties across five autonomous and persistent threads under a strict lock-synchronized schema to prevent race conditions and transaction deadlocks during continuous disk write operations.

To prevent filesystem I/O bottlenecks when logging dense network events, a collection of transaction-safe locks are instantiated and uniquely mapped to individual JSON lines (\texttt{.jsonl}) logging files. At the same time, independent memory-state locks are allocated to manage runtime mutations to shared memory containers,
preventing data corruption when multiple monitoring components concurrently access or modify relay data.

The framework restricts its data collection to all available relays on the network, a distinct target collection containing controlled static endpoints, a selected third-party endpoint probes, and a scalable pool of ephemeral addresses generated dynamically at runtime. By monitoring this mixed collection of hidden service 
categories, the system ensures comprehensive coverage across distinct regions of the HRT.

A five-thread colletion pipeline splits with each thread targeting a specific operational routine:

\begin{enumerate}
    \item \textbf{Statful Consensus Auditing and Sybil Vector Extraction:} These threads are driven by two independent consensus tracking routines, network-wide consensus monitor and tracked HSDirs status monitor. Triggered hourly at XX:05 UTC, these modules ingest the global and tracked directory network consensus information to map the state of active global and tracked relays. The audit engine reads the microdescriptor headers of every operational relay flagged with \texttt{HSDir, Stable} and \texttt{Fast} properties, mapping their IP addresses, nickname, open Onion Routing (OR) ports, and cryptographic fingerprints. Implemented with a microdescriptor parser bypass, sybil identity markers and flag clustered relays operated by coordinated entities are extracted with the capture of raw declared family lists and cryptographic \texttt{famil\_ids} vectors. Tracked active node that disappear from the consensus directory accumulates hourly absenses considering an offline churn; a consecutive absense count of $\geq$ 2 sets off an state reassignment to a dead-offline churn, capturing long-lived relay drop rates.
    \item \textbf{Active Outbound Prober:} A dedicated prober operating on a fixed periodic 20-minute loop. The mechanism randomizes the target request queue and dispatches raw asynchronous \texttt{HSFETCH} instruction lookups directly to the control port socket. These queries bypass Tor's internal circuit-level memory dependencies within the local Tor client, forcing the daemon to query teh remote directory nodes directly. This method discloses silent failues where a node listed as online cannot serve requested descriptors. To achieve complete coverage a two-tier lookup is triggered, a standard dynamic hash ring resolution, and a direct poll targeted at individual fingerprints historically recorded in memory as active hosting relays. These lookup patterns compels the local Tor daemon to query the remote relays directly, causing an asynchronous event listener to log the resulting network events.
	\item \textbf{Asynchronous Event Listener:} Thread instantiated to establish a persistent, isolated, direct subscription to the controller's internal \texttt{HS\_DESC} stream event pipeline, which tracks real-time descriptor uploads, fetch requests, and operational failures.
    \item \textbf{Scaled Ephemeral Manager:} To scale the observation scope and neutralize spatial placement bias inherent to a set of static domains, this runtime allocation script automatically invokes Tor control commands to dynamically deploy fifty detached , ephemeral hidden service endpoints and cycled at regular daily intervals. These generated cryptographic addresses are automatically appended to the active tracking set memory, and spreaded across the HRT ensuring broader statistical sampling across the decentralized DHT ring space.
\end{enumerate}

\subsection{Hash Ring Mapping and Low-Level Protocol Parsing}

A significant engineering barrier rised upon the degradation of the primary Python control library \texttt{Stem} for several years since June 2023, which lacks support of native parsing tailored to the modern version 3 onion service architecture. To bypass this limitation without sacrificing data granularity, 
custom protocol-level bypasses were introduced. Instead of using the broken library abstractions, the script builds isolated channels to prevent control port command congestion and interact directly with the raw socket stream. 

While the targeted consensus auditing can directly query (\texttt{.get\_microdescriptor()}) and partially parse relay consensus streams, the global consensus audit downloads unparsed consensus streams via raw \texttt{get\_info("dir/status-vote/current/consensus-micr} \texttt{odesc")}. It then extracts critical 
metrics and processes the text files using custom-built parsing that manually extracts the global SRV, decodes base64-encoded relay components, and handles raw-byte structures to compute the 32-byte HRT coordinate positions. 

These bypasses allows the framework to replicate the 32-byte DHT ring topology, which by capturing the global SRV and individual relay identity keys ($\textit{K}_{\text{ID}}$) from the consensus
, and parsing the current time period ($\textit{T}_{\text{p}}$), calculates the precise position index of every directory node using native SHA3-256 hashing routine:
\begin{equation}
\mathcal{R}_{\text{index}} = \text{SHA3-256}\big(\textit{K}_{\text{ID}} \parallel \textit{T}_{\text{p}} \parallel \textit{SRV} \big)
\end{equation}
The tracked consensus contrary to the network-wide consensus exclusively extracts the unique 40-character hexadecimal hash ring coordinate (\texttt{HSDIR\_INDEX}), the base64-encoded blinded public descriptor identifier which requires knowing the target 35-byte onion service address, and the server-side validation 
error reasons from the asynchronous event listener, which populates the shared memory container for targeted directory nodes. Then evaluates the clockwise distance $\Delta$ between a directory's coordinate $\textit{I}_{\text{HSDir}}$ on the HRT and the target onion service descriptor index $\textit{I}_{\text{onion}}$ 
across the circular ring space $\mathcal{M} = 2^{256}$:
\begin{equation}
\Delta = \left( I_{\text{HSDir}} - I_{\text{onion}} \right) \pmod{2^{256}}
\end{equation}
This distance is subsequently converted into a normalized position percentage relative to the total ring space to chart wether monitored directories are distributed uniformly or clustered around specific identities.

\subsection{Temporal Synchronization and Boundary Isolation}

A foundational requisite of the measurement system is the management of daily SRV transaction and the subsequent reshuffling of the hash ring coordinates and active blinded keys, which Tor V3 protocol enforces the rotation exactly at 00:00 UTC. 

To prevent data contamination from overlapping consensus epochs, the framework
enforces a strict exclusion boundary window spanning from 23:55 UTC to 00:07 UTC. During this interval, live network probes and live event captures are suspended or systematically filtered.

To provide directory authorities sufficient time to settle their consensus voting and distribute the updated SRV across the network, the framework uses staggered offsets to resume tracking:

Old ephemeral service nodes are terminated and newly deployed at 00:01 UTC. Precisely at 00:05 UTC, upon verifying a calendar day or month rollover, the active runtime directory nodes on 
shared memory, network-wide and targeted, are state verified on the new calendar rollover over a script routine, archived to a disk \texttt{.jsonl} log file, flushed from shared memory, and
re-initialized to a blank vector slate. Active outbound network probes resume at 00:08 UTC, and tracked and network-wide consensus audits re-engaged at 00:10 UTC. 

This synchronized time structure ensures that all local tracking states, rotated blinded keys, and newly designated remote directory authorities have fully stabilized across the network 
before data collection resumes, preserving high-fidelity of long-term churn measurements.

\subsection{Structural Sybil Footprint Indentification}

An analytical objective is differentiate random, isolated relay failures from systematic disruptions managed by large-scale coordinated node operators. When a monitored HSDir drops from the network
or fails to serve valid descriptor, the framework collects structural identity markers to detect potential Sybil concentrations based on selected criteria:

\begin{itemize}
    \item \textbf{Network Prefix Aggregation:} The framework maps IPv4 and IPv6 addresses of every directory node into standard \texttt{/16} and \texttt{/24} subnets. This grouping isolates cases where a single operator deploys multiple discrete directory nodes into the consensus within identical data center allocations or IP routing subnets to artificially influence a section of the HRT.
    \item \textbf{Declared and Effective Administrative Families:} The microdescriptor records are parsed hourly to extract both declared relay family identifiers and mutual recognized effective family lists. This tracks evaluate compliance with Tor's administrative policies and flags undeclared Sybil networks that exhibit identical uptime profiles or infrastructure features while omitting formal administrative declarations.
\end{itemize}

\begin{table}[htbp]
\caption{Table Type Styles}
\begin{center}
\begin{tabular}{|c|c|c|c|}
\hline
\textbf{Table}&\multicolumn{3}{|c|}{\textbf{Table Column Head}} \\
\cline{2-4} 
\textbf{Head} & \textbf{\textit{Table column subhead}}& \textbf{\textit{Subhead}}& \textbf{\textit{Subhead}} \\
\hline
copy& More table copy$^{\mathrm{a}}$& &  \\
\hline
\multicolumn{4}{l}{$^{\mathrm{a}}$Sample of a Table footnote.}
\end{tabular}
\label{tab1}
\end{center}
\end{table}

\begin{figure}[htbp]
\centerline{\includegraphics[width=0.7\textwidth]{images/logo_UC3M.png}}
\caption{Example of a figure caption.}\label{fig}
\end{figure}

\section*{Acknowledgment}

The preferred spelling of the word ``acknowledgment'' in America is without 
an ``e'' after the ``g''. Avoid the stilted expression ``one of us (R. B. 
G.) thanks $\ldots$''. Instead, try ``R. B. G. thanks$\ldots$''. Put sponsor 
acknowledgments in the unnumbered footnote on the first page.

\section*{References}

Please number citations consecutively within brackets \cite{lamp94}. The 
sentence punctuation follows the bracket. Refer simply to the reference 
number, as in \cite{lesk77}---do not use ``Ref. \cite{lesk77}'' or ``reference \cite{lesk77}'' except at 
the beginning of a sentence: ``Reference \cite{lesk77} was the first $\ldots$''

\printbibliography

\appendices

\section{DECLARATION OF USE OF GENERATIVE ARTIFICIAL INTELLIGENCE (AI)}
\label{app:ai_use}



\end{document}